\section{A lossless weak estimate}\label{sec:weak}

The power saving of \Cref{s:sparse} is available when the sparsity
parameter is a positive power of the packet scale.  We now reach every
$D\geq1$ by a second induction.  Its geometric step places most of a
hypothetical large level set near affine plates.  The frame estimate can
then be applied in one fewer spatial dimension.  The additional unit
density parameter in the next theorem makes the small-density part of
this induction contract.

\begin{theorem}[Weak estimate with density and sparsity]\label{w:weak}
Fix $n\geq3$, and use the compact arrays, measure conditions and fixed
time range of \Cref{d:measure}.  There are constants
\[
 p>2,\qquad \sigma>0,\qquad 0<2a<1-\beta,
 \qquad \beta=\frac2{n+1},\quad e=\frac{n^2}{n+1},
\]
depending only on the dimension and the fixed packet conventions, for
which the following holds.  Suppose that the array has coefficient
energy $E$, that $1\leq D\leq L$, and that the measure has ball and
packet-moment parameter $m$ as in \Cref{d:measure}.  Suppose also that
\[
 \mu(B(z,1))\leq m\chi,\qquad 0<\chi\leq1,
\]
and that its support is in a box of side $O(L^2)$ within the fixed time
range.  For every $g>0$,
\begin{equation}\label{w:eq-weak}
 \mu\{z:|U(z)|\geq \chi^a gL^{-e}\}
 \leq C_{\mathrm{ind}}mL^{2n}g^{-p}D^{-\sigma}E^{p/2}.
\end{equation}
The constants are independent of $L,D,m,\chi,g$ and the particular
array.  The assertion also holds for a subarray with its original
native profiles.
\end{theorem}

We prove the theorem with $E=1$; multiplying all coefficients by
$E^{1/2}$ gives the general statement.  A zero energy or a zero measure
parameter is harmless.  Every recursive application will use original
compact profiles, possibly after the unrotated cap rescaling already
specified in \Cref{d:transport}.  In particular, derivatives introduced
in a nonrecursive localization will never become extra hypotheses at
the next induction level.

The induction is on dyadic ranges of $L$.  For a small exponent
$\kappa>0$ to be chosen below, the inner parameter at a fixed range
of $L$ is the number of doublings of $D$ needed to reach $L^\kappa$.  Thus an application at smaller $L$ is an outer induction
step, while an application at the same $L$ and at least twice the
value of $D$ is an inner step.  Angular rescaling will use the outer
hypothesis; removing subarrays with smaller packet moments will use
the inner hypothesis.  Spatial support constants remain fixed:
partition a larger box into a bounded number of boxes of the stipulated
size, keep the compact packets meeting each, and sum their
$p/2$-powers of energy.  The energy overlap is bounded.  The time-range
constant is fixed to contain the packet supports.

\subsection{Order of choices and the easy ranges}\label{w:choices}

Choose
\begin{equation}\label{w:eq-ab}
 \beta<b<\frac2n,\qquad 2a=1-b+\delta,
 \qquad \zeta=1-\beta-2a=b-\beta-\delta>0,
\end{equation}
with $\delta>0$ small.  We will use
\[
 X=BD,\qquad R=K^J\asymp X^u,
\]
where $K$ is dyadic, $J$ is a fixed large integer, $u>0$ is small, and
$B$ is a fixed large constant.  Dyadic rounding changes only constants
depending on $J$.  We list the order of choices because both inductions
use strict gains.

\begin{enumerate}[label=\textup{(\roman*)}]
\item Fix $b$, then small $\delta$ and a range $2<p\leq p_{\max}$.
Choose $J$ large enough for \Cref{w:angular}.  That lemma gives a fixed
gain $R^{-c_2}$, with $c_2>0$ uniform as an exponent in this range of
$p$.  Its numerical constants need not be uniform as $p\downarrow2$.
\item Fix a small unit-density cutoff exponent $v_0>0$.  The positive
dimensional margin
\begin{equation}\label{w:eq-threshold-margin}
 n-\frac12-e=\frac{n-1}{2(n+1)}>0
\end{equation}
allows the power cutoff in the eventual frame application to be fixed
at this stage.  For this fixed cutoff, fix the normalized profile
bounds and the envelope and derivative orders for \Cref{f:frame},
including fixed reserves for its nonrecursive localization.  The
inverse-transversality, coarse-tuple and witness costs are then
$R^{C_f}$, with $C_f$ independent of all geometric
powers chosen below.  Take $u$ so small that $uC_f<c_s/4$, reducing it
for any other fixed frame cost.  Then choose $\sigma>0$ much smaller
than $uc_2$ and $c_s$.
\item Choose the finite powers of $X$ required for the inner sparsity
parameter, the wall tests, the plate count and width, the final ratio
$L/h$, and the tiny-density cutoff.  The exponent in the eventual
polynomial bound on $g$ is also determined by these choices.  Use
$p_{\max}$ in any upper estimate involving the still unfixed $p$.
\item Choose $\kappa>0$ small enough that, in the remaining range
$D<L^\kappa$, $L$ dominates every specified power of $X$ at large
scales for fixed $B$.  Also impose
$\kappa\sigma\ll v_0\zeta$.  Choose $p-2>0$ only afterwards,
sufficiently small relative to $\kappa$ and to the eventual level
exponent.  Then choose the small power losses and the finite derivative
orders needed to make absolute errors negligible.  Finally choose $B$,
the base scale, and $C_{\mathrm{ind}}$, in that order.
\end{enumerate}

The sparse estimate may require more derivatives of the original
profiles after the geometric powers have been chosen.  This later
regularity reserve does not change the fixed frame inputs in step
\textup{(ii)}, and hence does not change $C_f$.  Its multiplicative
constants are absorbed by the final base scale and induction constant.

By coefficient Cauchy--Schwarz and phase-space counting,
$\norm{U}_\infty\lesssim1$.  Hence a nonempty level set has
\begin{equation}\label{w:eq-gtrivial}
 g\lesssim L^e\chi^{-a}.
\end{equation}
For $D\geq L^\kappa$, Chebyshev and \Cref{s:sparse} give
\begin{equation}\label{w:eq-largeD}
 \mu\{|U|\geq\chi^agL^{-e}\}
 \leq C_\kappa mL^{2n}D^{-c_s}\chi^{\zeta}g^{-2}.
\end{equation}
The ratio of this bound to the desired right side, apart from its
constant, is at most
\[
 D^{-(c_s-\sigma)}\chi^{\zeta}g^{p-2}
 \lesssim L^{e(p-2)-\kappa(c_s-\sigma)}
       \chi^{\zeta-a(p-2)}.
\]
It is bounded if $p-2$ is small enough.  In the range
$D<L^\kappa$ and $\chi\leq L^{-v_0}$, \Cref{d:power-loss} gives
the same comparison with
\[
 L^{\epsilon+\kappa\sigma+e(p-2)}
 \chi^{\zeta-a(p-2)}.
\]
Choose $\epsilon$ small after $\kappa,p-2$; the positive density
power makes this bounded as well.  These estimates also handle bounded
scales uniformly as $\chi\downarrow0$ by enlarging the final constant.
If $g<D^{-1}$, the total mass bound $\mu(\R^{n+1})\lesssim mL^{2n}$
is enough, since $g^{-p}D^{-\sigma}\geq D^{p-\sigma}\geq1$.

It remains to consider
\begin{equation}\label{w:eq-active}
 \chi>L^{-v_0},\qquad D^{-1}\leq g\lesssim L^{e+av_0},
 \qquad D<L^\kappa.
\end{equation}
In this range the threshold $\lambda=\chi^agL^{-e}$ is bounded
below by a fixed negative power of $L$.  Every absolute approximation
error below can consequently be assigned a sufficiently large fixed
negative power of $L$.

\subsection{An angular induction with curved narrow sets}

Besides a neighborhood of an affine codimension-two plane, we must
exclude a small quadratic sublevel set inside a hyperplane.  The latter
condition will detect the second fundamental form of a polynomial wall.
All frequency bins in the following lemma are bins of the original
unrotated grid.  Selecting a bin means selecting its entire field.

\begin{lemma}[Angular induction]\label{w:angular}
Assume the outer induction hypothesis in \Cref{w:weak} and the active
range \eqref{w:eq-active}.  Fix a subarray of energy $E_0\leq1$.
In every cube of a fixed spacetime grid of side $R^4$, prescribe some
of its frequency bins of side $1/R$.  Suppose their centers are covered
by $O((1+\log X)^{C_0})$ sets, each of one of the following types:
\begin{enumerate}[label=\textup{(\alph*)}]
\item a strip of width $O(R^{-1})$ about an affine $(n-2)$-plane;
\item the $O(R^{-1})$-neighborhood of
\[
 \{v\in H\cap B: |Q_H(v)|\leq R^{-4}\},
\]
where $H$ is an affine hyperplane, $B$ is a fixed bounded region,
and $Q_H$ has coefficient norm one in orthonormal coordinates on $H$
with bounded origin.
\end{enumerate}
The cover and the prescribed bins may depend on the macro cube.  Let
$U_{\mathrm{sel}}$ denote the resulting prescribed-bin sum on that cube.
For any fixed $c>0$,
\begin{equation}\label{w:eq-angular}
 \mu\{|U_{\mathrm{sel}}|\geq c\lambda\}
 \leq C C_{\mathrm{ind}}R^{-c_2}
       mL^{2n}g^{-p}E_0^{p/2}+C_NmL^{-N}.
\end{equation}
Here $N$ can be any prescribed fixed exponent, given sufficiently many
initial profile derivatives.  The exponent $c_2>0$ depends on the
dimensional choices in \eqref{w:eq-ab}, not on $u$ or $C_0$.
\end{lemma}

\begin{proof}
We give the geometric iteration, the measure transport, and the
exponent calculation separately.

\paragraph{The narrow geometry.}
Assign every selected bin to one covering set.  In a set of type
\textup{(b)}, choose one approximating point for each assigned bin and
sort the sizes of $\abs{\nabla Q_H}$ there into dyadic classes $m_1$,
down to $1/R$; put all smaller gradients in the last class.  There are
$O(\log R)$ classes.  Each assignment concerns the whole bin, not a
subset of its packet labels.

Consider one class through the $J$ iterations of ratio $K$.  If $m_1$
is a sufficiently small constant, the normalization of $Q_H$, its
bounded coefficients, and its small value at an approximating point
force its quadratic part to have norm bounded below.  Some component
of $\nabla Q_H$ then has a linear part bounded below.  The condition
$\abs{\nabla Q_H}\lesssim m_1$ puts this class in a strip of width
$O(m_1+R^{-1})$ within the $O(R^{-1})$-neighborhood of $H$.
For parent inverse widths
\[
 s\leq \frac{c}{Km_1},
\]
this is a codimension-two strip at the next child precision.

For a dyadic class with a positive lower gradient bound, the other
regime is
\[
 \frac{CK}{m_1}\leq s\leq \frac{R}{K}.
\]
Choose an approximating point in the parent.  Taylor expansion around
it places all assigned descendants within
\begin{equation}\label{w:eq-quadraticwidth}
 C\bigl(m_1^{-1}(s^{-2}+R^{-4})+R^{-1}\bigr)
 \lesssim (Ks)^{-1}
\end{equation}
of the tangent affine $(n-2)$-plane to its level set in $H$.
The last, small-gradient class never needs this second regime.  If the
gradient is bounded below by a fixed constant, only a bounded number
of initial steps are outside the second regime.  In general the gap
between the two regimes is a factor $O(K^2)$, so at most a bounded
number of ratio-$K$ steps need a crude estimate.

Use the functional $\mathcal T_b$ of \Cref{d:transport}, with
$P=2/(1-b)$.  Since $b<2/n$, we have $P\leq q_{n-2}$.
At every good step \Cref{d:strip-decoupling} applies after rescaling the width in
\eqref{w:eq-quadraticwidth} to $O(K^{-1})$.  The exceptional steps use
Cauchy--Schwarz.  Consequently the whole angular iteration, excluding
the transport factors, costs
\begin{equation}\label{w:eq-Kcost}
 C_\epsilon K^C R^\epsilon,
\end{equation}
where the exponent $C$ is independent of $J$.  Constants depending on
the fixed $J$ are allowed.  This distinction is essential: good-step
losses multiply to $R^\epsilon$, rather than $K^{CJ}$.

For this nonrecursive calculation one may smoothly truncate compact
profiles in native spacetime Fourier variables at radius $L^{1/4}$.
With sufficiently many fixed derivatives, every relevant subfield and
rescaling has error smaller than any prescribed negative power of $L$.
There are only polynomially many labels of packets meeting the domain.
The choice of $\kappa$ makes the fixed hierarchy $R=K^J$ small enough
relative to $L$ that the truncated supports satisfy the thickness and
overlap hypotheses of the packet estimates at each step.  Supremum
estimates for these approximations use ordinary convolution.  At the
terminal step we return to the original compact child fields, so none
of this truncation changes the induction's profile class.  The absolute
error estimate is justified below, including when the tested mass is
small.

\paragraph{Initial measure functional and recombination.}
Let $\nu$ be the original measure restricted to the set on the left
of \eqref{w:eq-angular}, and write $M_0=\nu(\R^{n+1})$.
For each cover/class index, pad the enumeration where needed and
restrict $\nu$ separately to the macro cubes.  On a unit cube $q$,
the density bound gives $\mu(q)^{1-b}\lesssim(m\chi)^{1-b}$.
H\"older on the constituent unit cubes, followed by the weighted
supremum estimate for the bandlimited approximation, bounds
$\lambda^2M_0$ by the sum of the initial $\mathcal T_b$ functionals,
multiplied by
\[
 C(m\chi)^{1-b}(1+\log X)^C.
\]
The sum of unit-cube convolution weights inside a $K^2$-cube is bounded
by the enlarged Schwartz weight used for that cube.  Thus the weighted
versions of the transport inequality apply without a factor counting
unit cubes.  During each calculation the angular subset prescribed on
the original macro is kept fixed, including in every shifted norm.

After $J$ steps, the transport factors and \eqref{w:eq-Kcost} give
\begin{equation}\label{w:eq-angular-factor}
 C_\epsilon K^C R^{\epsilon+2-(n+2)b}.
\end{equation}
We explain why the macro restrictions can now be recombined.  Write
\[
 A_{K,v}(x,t)=((x-2vt)/K,t/K^2).
\]
For a fixed angular chain, class index, and complete grid-shift
sequence, every step is a translated pushforward under $A_{K,v_j}$.
The grids on the right of the transport inequalities are global and
unrotated; rotations used to prove a strip inequality do not change
these lists.  Hence the composition is
\begin{equation}\label{w:eq-affine-composition}
 \tau_{z_{\rm sh}}\circ A_{R,V},\qquad
 V=v_1+v_2/K+\cdots+v_J/K^{J-1},\qquad |V|\leq C_n,
\end{equation}
with a translation $z_{\rm sh}$ independent of the original macro cube.
The inverse image of a terminal $K^2$-cube has diameter at most
$C_nR^2K^2\leq C_nR^4$, regardless of the size of $z_{\rm sh}$.  It therefore
meets only boundedly many original macros.  For their transported
measures this gives, on every terminal cube $Q$,
\[
 \sum_{\mathcal M}\nu_{\mathcal M}(Q)^b
 \leq C_n\left(\sum_{\mathcal M}\nu_{\mathcal M}(Q)\right)^b.
\]
The combined measure is dominated by one translated pushforward of
$\nu$, has mass at most $M_0$, and has ball parameter
\begin{equation}\label{w:eq-childm}
 m'\asymp C mR^{n+1}.
\end{equation}
Infinite shift sequences are summed with their summable weights;
neither the mass, the ball constant, nor the number of occupied cubes
depends on their absolute translation.  Since assignments were by
whole original $1/R$-bins, the terminal field is the same full compact
bin field, or zero, on every contributing macro.  Thus the outer
induction will charge each original bin energy once.

\paragraph{The terminal induction.}
Put $l=L/R$.  In each terminal $K^2$-cube with nonzero original child
supremum, move its mass to a point arbitrarily close to that supremum
and bound its $P$-norm by the supremum.  This changes ball bounds by
at most $K^C$.  Sort the cube masses as $\nu_{\mathrm{child}}(Q)\asymp
m'd$.  Unit covering of the inverse image, whose volume scales as
$R^{n+2}$ while \eqref{w:eq-childm} scales as $R^{n+1}$, gives
\begin{equation}\label{w:eq-dcap}
 d\leq CK^C R\chi.
\end{equation}
The moved measure has ball parameter $CK^Cm'$, relative unit density
at most $\min(1,Cd)$, and, after fixed constant inflation, admissible
moment parameter $D=1$.  Its relevant time range is that of the
original compact child field.  Spatial boxes can be partitioned as
in the theorem.

Apply the outer induction at $l$ and integrate the weak $L^p$ bound
against the trivial mass bound $M_0$.  Since $p>2$, the resulting
$L^2$ integral is bounded by
\begin{equation}\label{w:eq-terminalL2}
 CK^C C_{\mathrm{ind}}^{2/p}(m')^{2/p}
 l^{4n/p-2e}d^{2a}M_0^{1-2/p}E_{\mathrm{bin}},
\end{equation}
where $E_{\mathrm{bin}}$ is that bin's original coefficient energy.
Indeed, integrating a distribution estimate
$\nu\{|V|>t\}\leq\min(M_0,At^{-p})$ gives
$\int|V|^2\,\dd\nu\leq C_p A^{2/p}M_0^{1-2/p}$.
Multiplication by $(m'd)^{b-1}$ returns to the mass-power functional;
the fixed $K$-volume factor costs another $K^C$.  The dyadic $d$ sum
converges at its lower end because
\[
 b-1+2a=\delta>0.
\]
By \eqref{w:eq-dcap}, its total cost is
$C(K^CR\chi)^\delta$.

To justify returning from the Fourier approximations even on cubes
with zero original supremum, normalize $m=1$.  There are polynomially
many occupied cubes, and concavity gives
\[
 \sum_Q\nu(Q)^b\leq(\#Q)^{1-b}M_0^b.
\]
Choosing the fixed truncation order large enough makes the entire
error, with all prefactors restored, at most $L^{-N'}M_0^b$ for any
prescribed $N'$.  Either this absorbs into $\lambda^2M_0$, or, using
$\lambda\geq L^{-C}$ from \eqref{w:eq-active},
\[
 \lambda^2M_0\lesssim L^{-N'}M_0^b
 \quad\Longrightarrow\quad
 M_0\lesssim L^{-(N'-2C)/(1-b)}.
\]
The latter is the asserted absolute negligible mass after choosing
$N'$ large.  Restoring homogeneity supplies the factor $m$.

\paragraph{The gain.}
Sum the disjoint bin energies.  The factors of $m$ and $\chi$ before
solving for $M_0$ are exactly $m^{2/p}\chi^{2a}$.  Apart from
$K^CR^\epsilon$ and fixed logarithmic powers, the exponent of $R$ is
\begin{equation}\label{w:eq-Gamma}
 \Gamma(p,\delta)=2-(n+2)b+(n+1)(b-1+2/p)-4n/p+2e+\delta.
\end{equation}
At $p=2$, $\delta=0$, direct simplification gives
$\Gamma=\beta-b<0$.  Solving for $M_0$ multiplies this exponent by
$p/2>0$.  First choose $\delta$ and $p_{\max}-2$ small, then $J$
large so that $K^C=R^{C/J}$ is paid by the strict margin, and finally
choose the decoupling loss small.  The remaining logarithmic factors
are absorbed by increasing $B$ for the fixed $C_0$.  This proves
\eqref{w:eq-angular} with a fixed $c_2>0$.
\end{proof}

\subsection{Heavy bins and a polynomial wall}\label{w:heavy}

Suppose, towards a contradiction to \Cref{w:weak}, that
\begin{equation}\label{w:eq-counterlevel}
 Y=\{|U|\geq\lambda\},\qquad
 M_Y=\mu(Y)>C_{\mathrm{ind}}mL^{2n}D^{-\sigma}g^{-p}.
\end{equation}
All constants used in the geometric construction below are independent
of $C_{\mathrm{ind}}$.  Translations of the spatial box are immaterial.
For a macro cube with center $z=(x,t)$, call a $1/R$-bin heavy if
\begin{equation}\label{w:eq-heavy}
 \sum_{\alpha\ \mathrm{in\ the\ bin}}|c_\alpha|^2
 \left(1+\frac{|x-y_\alpha-2\omega_\alpha t|}{L}\right)^{-M}
 \geq W,\qquad
 W=X^{-P_0}\frac{mL^{n+1}}{M_Y}.
\end{equation}
The power $P_0$ will be chosen after an inner sparsity exponent $P_1$.

\begin{lemma}[Reduction to broad heavy incidences]\label{w:heavy-reduction}
There is a subset $Y'\subset Y$ with $\mu(Y')\geq cM_Y$ such that,
at each point of $Y'$:
\begin{enumerate}[label=\textup{(\roman*)}]
\item the heavy-bin field has modulus at least $\lambda/2$;
\item the heavy-bin centers are not contained in a strip of any fixed
large constant times $R^{-1}$ about an affine $(n-2)$-plane, nor in a
single quadratic neighborhood of the form in \Cref{w:angular};
\item every heavy bin supplies at least $W/2$ of its weighted
coefficient square mass from trajectories within distance $X^{C_w}L$.
\end{enumerate}
In particular some $n$ heavy bins have direction vectors
$(2\omega,1)$ whose $n$-volume is at least $R^{-C}$ throughout the
bins.
\end{lemma}

\begin{proof}
For each bin restrict $\mu|_Y$ to the macros where that bin is light.
Since $R^4\ll L$, the weight in \eqref{w:eq-heavy} is comparable at
the macro center and at each point of the macro.  Therefore the sum of
its coefficient energies times their packet moments on this restricted
measure is at most $CWM_Y$.

Put $D_1=X^{P_1}$.  Split off the labels whose moment exceeds
$mL^{n+1}/D_1$.  In each bin their total energy is at most
$CD_1X^{-P_0}$; summing over the $O(R^n)$ bins only contributes a fixed
power of $R$.  The other labels satisfy the inner induction hypothesis
at $D_1$, on the restricted measure, with the unchanged original
profiles.  Apply it at threshold $cR^{-n}\lambda$.  Disjoint bin
energies and $p/2>1$ give total exceptional mass at most
\begin{equation}\label{w:eq-light-inner}
 C C_{\mathrm{ind}}mL^{2n}g^{-p}R^{np}D_1^{-\sigma}.
\end{equation}
Choose $P_1$ so that $2D\leq D_1$ and
$un p_{\max}-\sigma P_1+\sigma<0$ with a fixed spare margin.
The later choice of $\kappa$ ensures $D_1\leq L$.  If $D_1$ has
already reached $L^\kappa$, the large-$D$ case is available; otherwise
this is a strict inner induction step.

For each split-off bin, use unrotated cap scaling and the outer
induction at $L/R$.  The ball parameter is $CmR^{n+1}$, the relative
unit density is at most $\min(1,C\chi R)$, and the scaled level
parameter is at least $gR^{-C}$.  Consequently all split-off fields
together contribute at most
\begin{equation}\label{w:eq-light-outer}
 C C_{\mathrm{ind}}mL^{2n}g^{-p}
 R^C(D_1X^{-P_0})^{p/2}.
\end{equation}
The fixed power $R^C$ includes the bin count and threshold splitting;
its exponent can use $p_{\max}$.  Choose $P_0$ sufficiently larger
than $P_1$ so that this expression, and \eqref{w:eq-light-inner}, are
small fractions of the lower bound in \eqref{w:eq-counterlevel} once
$B$ is large.  The thresholds of the $O(R^n)$ light-bin fields sum to
at most $\lambda/2$ when $c$ is small.  This proves \textup{(i)}.

On every adverse macro for \textup{(ii)}, the heavy-bin sum is exactly
a prescribed-bin field of the type in \Cref{w:angular}.  Its mass is
at most the target bound times $CR^{-c_2}D^\sigma$, plus an absolute
negligible error.  Since $\sigma<uc_2$, increasing $B$ makes this a
small fraction.  The construction may choose the adverse hyperplane
and quadratic separately on each macro, as allowed by that lemma.

For \textup{(iii)}, the measure ball bound and the fixed time range
give, by covering tube annuli,
\[
 \int_{|x-y_\alpha-2\omega_\alpha t|>X^{C_w}L}
 \left(1+\frac{|x-y_\alpha-2\omega_\alpha t|}{L}\right)^{-M}
 \,\dd\mu\leq CmL^{n+1}X^{-A_w},
\]
where $A_w$ can exceed $P_0$ by choosing $C_w$ large (the fixed
exponent $M$ was chosen sufficiently large first).  Sum with
$|c_\alpha|^2$ and compare with $WM_Y=X^{-P_0}mL^{n+1}$.
Markov's inequality removes only a small fraction of $Y$ and ensures
the stated contribution in every heavy bin.  Finally, successive
maximal affine heights among the remaining centers give $n$ robustly
independent directions; failure would place all centers in a forbidden
codimension-two strip.  Choosing the strip constant large makes the
volume bound stable throughout the bins.
\end{proof}

Pass to tube coordinates $(x,t)/L$ and replace the measure by
$\nu=\mu/(mL^n)$ in those coordinates.  The relevant domain has diameter
$O(L)$ and
\[
 \nu(B(z,r))\leq Cr^n\quad(r\geq1).
\]
Write $N_Y=\nu(Y)$.  Formula \eqref{w:eq-heavy} becomes
\begin{equation}\label{w:eq-Wnormalized}
 W=X^{-P_0}L/N_Y.
\end{equation}
The incidence tubes have width $X^{C_w}$ in these coordinates.

\begin{lemma}[Incidence mass and a wall]\label{w:wall}
The preceding level set satisfies
\begin{equation}\label{w:eq-incidence-mass}
 N_YW^{n/(n-1)}\leq X^C,
 \qquad N_Y\geq X^{-C}L^n.
\end{equation}
Moreover, a subset of $Y'$ of mass comparable to $N_Y$ lies within
distance $X^C$ of the zero set of a polynomial of degree at most
$X^{C_{\mathrm{part}}}$ in tube coordinates.
\end{lemma}

\begin{proof}
For a robust $n$-tuple, project its directions onto an $n$-dimensional
space where their determinant is at least $R^{-C}$.  Such a space
exists by their $n$-volume lower bound.  A further polynomially fine
angular subdivision makes the projection common to the tuple's
families.  Apply endpoint $n$-linear Kakeya to the projected tubes.

To relate its cube sum to $\nu$, divide into columns of side comparable
to the tube width.  Each tube meets only polynomially many width boxes
over a projected column, uniformly in its label; its projected
direction is bounded away from zero by the determinant condition.
Each box has $\nu$-mass at most $X^C$.  If $a_{i,j}\geq0$ denotes the
incident weight of family $i$ in the $j$th vertical box, H\"older with
$n$ factors followed by the monotonicity of finite sequence norms gives
\[
 \sum_j\prod_{i=1}^n a_{i,j}^{1/(n-1)}
 \leq\prod_{i=1}^n\left(\sum_j a_{i,j}\right)^{1/(n-1)}.
\]
The total coefficient weight in each family is at most one.  The
projected Kakeya bound, its polynomial transversality cost and the
polynomial number of tuples therefore bound the mass seeing weight
$W$ from every family by $X^CW^{-n/(n-1)}$.  This proves the first
inequality in \eqref{w:eq-incidence-mass}.  Substitute
\eqref{w:eq-Wnormalized} to obtain the second.  We also have
$N_Y\lesssim L^n$ from the ball bound.

The separation into cells and a polynomial wall follows the
Fourier-restriction partitioning framework of
Guth \citep[Sections~3.2--3.4]{Guth2016Restriction}.
For completeness, use polynomial partitioning with parameter
$P_2=X^{C'}$ on $\nu|_Y$.  It gives a polynomial of degree $O(P_2)$,
$O(P_2^{n+1})$ sign regions, and mass at most
$CN_Y/P_2^{n+1}$ in each open region.  The finite-set partitioning
theorem follows from iterative polynomial bisection
\citep[Theorem~4.1]{GuthKatz2015}; its integrable-density form is
\citet[Theorem~1.4]{Guth2016Restriction}.  For the present finite Borel
measure, smooth it first and use that bisection argument.  Normalize the
coefficient vector of each of the finitely many factors separately
and take a convergent subsequence.  A compact subset off the limiting
zero set has stable signs; its limiting mass obeys the same bound.
Exhausting an open sign region by such compact subsets proves the
assertion for the original measure.  Mass on the zero set is already
part of the wall.  Connectedness of sign regions is not required.

A central line not contained in the wall meets only $O(P_2)$ open
sign regions.  Let $E_{\mathcal O}$ be the total coefficient energy of
labels whose line visits $\mathcal O$.  Then
$\sum_{\mathcal O}E_{\mathcal O}\lesssim P_2$.  Regions with
$E_{\mathcal O}>C_*P_2^{-n}$ have total mass at most $CN_Y/C_*$ by
the region mass bound.  Choose $C_*$ large.  Away from a width-enlarged
wall, a point and its nearby incident line point are in the same
region.  In every remaining region the preceding projected Kakeya
argument with local energies bounds the off-wall mass of $Y'$ by
\[
 X^CW^{-n/(n-1)}P_2^{-n^2/(n-1)}
 \leq X^{C''}N_YP_2^{-n^2/(n-1)}.
\]
The last inequality uses \eqref{w:eq-Wnormalized} and
$N_Y\lesssim L^n$.  Since
\[
 \frac{n^2}{n-1}-(n+1)=\frac1{n-1}>0,
\]
choosing $C'$ large pays the number of regions.  A fixed positive
fraction of the mass must therefore be near the polynomial wall.
\end{proof}

\subsection{Why the wall is flat}\label{w:curvature}

Testing the second fundamental form on incident line directions is
related to the flat-point argument of
Guth--Katz \citep[Lemma~3.3 and Corollary~3.4]{GuthKatz2010Algebraic}.
Here the incidences are approximate and the ambient dimension is higher;
the quadratic narrow estimate supplies the additional directional
exclusion needed to bound the full Hessian.  We prove the quantitative
reduction to affine plates below.

We now work in macroscopic coordinates, obtained by dividing tube
coordinates by $L$.  The domain is a fixed bounded box, and the wall
and incidence errors are at most
\begin{equation}\label{w:eq-Delta}
 \Delta\leq X^{C_4}/L.
\end{equation}
The exponent $C_4$ is already fixed.  All further enlargements in this
subsection are fixed powers of $X$ and will be small on the
macroscopic scale by the later choice of $\kappa$.

We use the following two elementary consequences of fixed-dimensional
semialgebraic complexity.  Here a description of complexity $X^C$
means that the number and degrees of its defining polynomials are
bounded by such a power; their coefficients are unrestricted.
\begin{enumerate}[label=\textup{(\alph*)}]
\item A bounded set of dimension at most $n-1$ and complexity $X^C$
is covered by $X^{C'}\rho^{-(n-1)}$ boxes of diameter $O(\rho)$.
To see this, use a generic shifted mesh.  A box meeting a component
that does not meet any face contains that component, so the component
bound pays such boxes.  Intersections with grid faces generically lower
dimension, have the same type of complexity bound, and are counted by
induction on dimension, with the zero-dimensional component bound as
base case.
\item A bounded set of complexity $X^C$ has a partition into $X^{C'}$
cells within which any two points can be joined by a path of length
$X^{C'}$.  Use a compatible cylindrical decomposition, including the
enclosing box.  Recursively lift paths in the base to a graph or to
the continuous middle section of a sector, adding vertical segments
inside a sector.  The paths have polynomial description complexity
in fixed dimension.  Each bounded coordinate has polynomially many
monotonicity pieces, so the length is bounded by their total variation.
The paths may be taken continuous and piecewise smooth.
\end{enumerate}
The required cylindrical decomposition, projection and bounded-variable
quantifier-elimination bounds are standard; see \citet{Basu2014}.
These bounds are polynomial in description size when dimension is
fixed and do not depend on coefficient magnitudes.  They control the
number of pieces and path lengths, not graph derivatives.  The latter
are obtained next.

\begin{lemma}[Regular graph pieces and derivative bounds]\label{w:regular-wall}
Let $Z$ be the square-free polynomial wall.  It has a semialgebraic
exceptional subset of dimension at most $n-1$ and complexity $X^C$
such that, outside its $C\rho$-neighborhood, every relevant point of
$Z$ has a local graph coordinate with bounded slope and second and
third graph derivatives bounded by a fixed power of $\rho^{-1}$.
The estimates are independent of the polynomial coefficients.
The mass of the discarded neighborhood in tube coordinates is at most
$X^C\rho L^n$.
\end{lemma}

\begin{proof}
Square-freeness implies that the singular set has dimension at most
$n-1$.  Decompose the bounded regular part into polynomially many
pieces with a specified omitted coordinate $j$, on which the implicit
graph $x_j=\phi(\bar x)$ has slope bounded by a dimensional constant.
Choose a finite set of base directions sufficient to polarize tensors
through order three.  Refine the pieces by the signs, including zero,
of the corresponding fourth directional derivatives of $\phi$.
Implicit differentiation makes these tests semialgebraic with
polynomial complexity.  Put all singular points, lower-dimensional
pieces and frontiers into the exceptional set; use a strictly larger
box so that its artificial boundary is harmless.  The frontier
dimension property gives the stated dimension bound.

At distance $C\rho$ from this set, a straight base segment of length
$c\rho$ lifts in the same graph/sign piece.  Indeed its lift has
bounded speed.  If maximal continuation stopped earlier, its Cauchy
endpoint would be a frontier or a singular point within distance
$\rho$, unless the implicit graph could be continued.  Both alternatives
contradict the stopping assumption.  Global injectivity of the base
projection is unnecessary.

On such a two-sided segment let $g(t)=\phi(u+te)$, so $|g'|\leq C$.
The fixed sign of $g''''$ makes $g'''$ monotone.  If
$|g'''(0)|=A$, on one half-interval of length comparable to $\rho$ its
magnitude is at least $A$, with constant sign.  The midpoint second
difference of $g'$ on that half-interval has magnitude at least
$cA\rho^2$ and at most $4C$.  Thus $A\lesssim\rho^{-2}$.
Apply the same argument on shorter interior intervals, then a
difference quotient for $g'$, to get $|g''|\lesssim\rho^{-1}$.
Polarization bounds the full second and third tensors.

By property \textup{(a)}, the exceptional set is covered by
$X^C\rho^{-(n-1)}$ boxes.  Each becomes a ball of radius $O(\rho L)$
in tube coordinates and has $\nu$-mass at most $C(\rho L)^n$.
Their total mass is $X^C\rho L^n$.  Choose
$\rho=X^{-C_5}$ with $C_5$ large and require $\Delta\ll\rho$.
By \eqref{w:eq-incidence-mass}, this discards an arbitrarily small
fraction of $N_Y$.
\end{proof}

Every remaining point of the level set near the wall can now be
attached to a regular $p\in Z$ at distance $O(\Delta)$, with the
preceding slope and derivative bounds.  No regularity of this
attachment is needed; all bad events below can be described by
existence of an eligible nearby $p$.

\begin{lemma}[Line tests]\label{w:line-tests}
For any prescribed large $A$, choose a fixed $C_6$ sufficiently large.
For every label, outside polynomially many macroscopic time intervals
of total length at most $X^{-A}$, an incidence with any eligible $p$
satisfies
\begin{equation}\label{w:eq-line-tests}
 |v_j-\nabla\phi(p)\cdot\bar v|\leq X^{C_6}/L,
 \qquad |\nabla^2\phi(p)[\bar v,\bar v]|\leq X^{C_6}/L,
 \qquad v=(2\omega,1).
\end{equation}
The exponent bounding the number of intervals is independent of the
sizes of $A,C_6$.
\end{lemma}

\begin{proof}
Lift the projection of the nearby line position onto the local graph.
The lift is $O(\Delta)$ from $p$, and its first and second derivatives
differ from those at $p$ by $O(\Delta X^C)$ by
\Cref{w:regular-wall}.  For the macroscopic line $\ell(s)$ define
\[
 H(s)=\ell_j(s)-\phi(\bar\ell(s)).
\]
At a relevant incidence, $|H|\lesssim\Delta$.
Failure of \eqref{w:eq-line-tests} forces $|H'|$ or $|H''|$ to be at
least $cX^{C_6}/L$, after increasing $C_6$ beyond the fixed derivative
error powers.

Order the regular roots over the projected line parameter by
cylindrical decomposition and refine by these inequalities and signs.
This gives polynomially many smooth branch intervals.  They are graph
sections, not sectors, and their implicit derivatives are the actual
derivatives of $H$.  On an interval where $|H|\lesssim\Delta$ and a
first derivative of fixed sign has magnitude at least $q$, the length
is $O(\Delta/q)$.  For a second derivative of fixed sign and magnitude
at least $q$, a midpoint second-difference estimate bounds the length
by $O((\Delta/q)^{1/2})$.  Set $q=cX^{C_6}/L$ and use
\eqref{w:eq-Delta}.  Increasing $C_6$ makes the sum of lengths at most
$X^{-A}$.  Thresholds change coefficients, not degrees or the number
of quantified variables, so the interval-count exponent does not
increase with $A,C_6$.  Quantifying the existence of any eligible
nearby point preserves this complexity bound and proves the last
assertion as well.
\end{proof}

The bad parameters of one label contribute at most
\begin{equation}\label{w:eq-badline-mass}
 X^C(LX^{-A}+X^C)
\end{equation}
to its enlarged tube's $\nu$-mass: cover the bad intervals in tube
coordinates by width boxes and use the ball bound.  The prefactor and
count exponents are fixed independently of $A,C_6$; the initial
$\Delta$ can be enlarged to include the incidence width.
Choose $A$ large, then $C_6$, and finally let $L$ dominate the fixed
powers.  Sum \eqref{w:eq-badline-mass} with $|c_\alpha|^2$ and compare
with $WN_Y=LX^{-P_0}$.  Markov removes only a small fraction of $Y'$,
and ensures that each heavy bin supplies at least one label passing
both tests.  In fact it supplies positive weighted mass of such
labels.  All assertions are simultaneous at each retained wall point.

\begin{proposition}[Affine plate covering]\label{w:geometry}
A subset $Y_0\subset Y'$ with $\mu(Y_0)\gtrsim M_Y$ is contained
in a union of at most $X^C$ affine hyperplane plates, each of thickness
$X^C$ in tube coordinates.
\end{proposition}

\begin{proof}
At a retained regular point $p$, let $\mathbf n=(\mathbf n_x,
\mathbf n_t)$ be a unit normal.  The first line test gives
\[
 |\mathbf n_x\cdot2\omega+\mathbf n_t|\lesssim X^{C_6}/L
\]
for a supplied bounded frequency, so $|\mathbf n_x|\gtrsim1$.
The exact tangent-frequency hyperplane is consequently well conditioned.

Suppose the graph Hessian has norm greater than $X^{C_7}/L$.
Normalize its quadratic form and restrict it via
$\omega\mapsto\overline{(2\omega,1)}$ to that exact hyperplane.
The resulting polynomial on the hyperplane has coefficient norm
comparable to one.  To verify this, choose coordinates on the tangent
space in which its time-one slice has bounded offset.  Its direction
space and one bounded point span the tangent space with bounded
condition number because $|\mathbf n_x|\gtrsim1$.  A homogeneous
quadratic has the form
\[
 Q(y,t)=y^TCy+2t b\cdot y+at^2;
\]
its restriction to $t=1$ records all of $C,b,a$.  The bounded graph
slope makes base projection another bounded isomorphism.  This proves
the coefficient comparison.  Divide the restricted polynomial by its
coefficient norm, so that it has norm exactly one as required in
\Cref{w:angular}; this changes the following bounds only by fixed factors.

Move each supplied frequency to the exact tangent hyperplane.  Its
distance is $O(X^{C_6}/L)$, and by the second line test its normalized
quadratic value there is at most
\[
 C X^{C_6-C_7}+C X^{C_6}/L.
\]
Choose $C_7$ large relative to $C_6,u$ and then make $L$ dominant.
This quantity is less than $R^{-4}$.  Every heavy-bin center is within
$O(R^{-1})$ of these frequencies, so all heavy bins are trapped in a
forbidden quadratic neighborhood from \Cref{w:heavy-reduction}.
The contradiction proves
\begin{equation}\label{w:eq-smallhess}
 \norm{\nabla^2\phi(p)}\leq X^{C_7}/L.
\end{equation}
This is where the quadratic narrow estimate is essential: tangent
directions on a nonflat ruled quadric may span an affine hyperplane,
but their quadratic values still detect its curvature.

For each fixed graph coordinate $j$, partition the exact semialgebraic
subset defined by regularity, the bounded slope condition and
\eqref{w:eq-smallhess} into cells with property \textup{(b)} above.
Use paths inside these cells, not paths inside the larger regular
wall.  Even if the base projection has several sheets, the implicit
gradient
\[
 g_j(p)=-\partial_{\mathrm{base}}P(p)/\partial_jP(p)
\]
is a single-valued function of the wall point.  Along a piecewise
smooth path $p$ in this subset, with base coordinate $u$,
\[
 \dd g_j=\nabla^2\phi_j\,\dd u,\qquad
 \dd p_j=g_j\cdot\dd u.
\]
For path length $D_0\leq X^C$, these identities bound the gradient
change by $X^{C_7}D_0/L$ and the height error from the starting tangent
plane by $X^{C_7}D_0^2/L$.  A projected loop reaching a different
sheet obeys the same identities.  Thus every cell lies within
$X^C/L$ of one affine plane in macroscopic coordinates.  Add the
$O(\Delta)$ attachment error and return to tube coordinates.  The
number of cells and the plate widths are both fixed powers of $X$,
as claimed.
\end{proof}

\subsection{Assigning packets to plates}\label{w:assignment}

The plate covering is used only to organize the original compact
packets.  It does not change their profiles.  Assign a label to a plate
if its central line comes within distance $X^{C_8}$ of that plate
somewhere in the enlarged tube-coordinate domain and its normal
velocity deviation for the same plate is at most $X^{C_9}/L$.
If there are several choices, choose one.  Thus the assigned coefficient
energies are disjoint.  Increase each assigned plate's width by a fixed
power of $X$ so that it contains the compact packet support in the
domain: the normal displacement over time $O(L)$ is at most $X^{C_9}$,
in addition to the initial distance and the fixed packet width.

Every unassigned label has moment on $\mu|_{Y_0}$ at most
$mL^{n+1}/D_1$, after choosing the powers in order.  Indeed choose
$C_8$ larger than the plate-width power and truncate its moment at
distance $X^{C_8}/2$, paying the tail by the same annular estimate as
in \Cref{w:heavy-reduction}.  If its line is ever this close to a
plate but is unassigned, its normal velocity exceeds $X^{C_9}/L$.
The time spent within that distance is at most
\[
 CL\frac{X^{C_8}+X^{C_{\mathrm{plate}}}}{X^{C_9}}
\]
for each plate, in tube coordinates.  Cover those intervals by width
boxes and sum over the polynomial number of plates.  Choosing $C_9$
large makes the resulting moment smaller than the asserted bound.
Inner induction at $D_1$ therefore removes the unassigned field on
only a small fraction of $Y_0$.  The assigned sum is at least a fixed
multiple of $\lambda$ on a set of mass comparable to $M_Y$.

Choose a dyadic inverse cap width $h$, nested with $R$, such that
\begin{equation}\label{w:eq-h}
 b_1=L/h\asymp X^{C_{10}}.
\end{equation}
The exponent $C_{10}$ is larger than the preceding plate-count and
width powers and the loss in \eqref{w:eq-incidence-mass}.  Intersections
of two enlarged plates whose projective normals differ by at least
$1/h$ have small total mass.  In tube coordinates such an intersection
fits in a box with $n-1$ lengths $O(L)$ and two remaining widths at
most $X^Ch$.  Covering by balls of the latter radius gives mass at most
$X^CL^{n-1}h$.  The same bound, with a larger fixed power, pays all
pairs of plates.  By \eqref{w:eq-incidence-mass} and
\eqref{w:eq-h}, this is negligible relative to $N_Y$.

Group the plates by a bounded-overlap grid of projective normals at
resolution $1/h$, choosing labels disjointly.  At a retained point only
boundedly many clusters contribute, so one cluster field has modulus
$\gtrsim\lambda$.  Cluster energies partition a subenergy of the
original array.  Within each cluster all frequencies lie within $C/h$
of a single affine frequency hyperplane.  To see this, the assignment
condition gives an almost tangent direction of bounded velocity for
each plate, so the spatial part of its normal is bounded below; changing
the normal by $O(1/h)$ changes its tangent-frequency hyperplane by the
same order.  A bounded spatial rotation, shear and modulation now put
this hyperplane in tangential frequency coordinates of dimension
$k=n-1$.  These transformations are used only in the following
nonrecursive estimates.

For each cluster, subdivide labels in the original frequency axes into
$1/R$-bins and then $1/h$-bins.  On a macro cube where the cluster is
high, call a coarse bin significant if its field has a witness there
of modulus at least $cR^{-n}\lambda$.  The nonsignificant fields sum
to at most a small multiple of $\lambda$ at a tested point.

Except for a small total mass, the tangential projections of the
significant centers cannot be covered by $O(C_{11}\log X)$ strips of
width $C_*/R$ about affine $(k-1)$-planes.  Here $C_*$ is any sufficiently
large fixed dimensional constant, and $C_{11}$ is the later fixed band
cutoff exponent.  Such a cover, together with the cluster's $C/h$
normal width, would put the original bins in codimension-two strips
in full frequency space.  Apply \Cref{w:angular} to the cluster subarray
on each adverse macro.  Summing disjoint energies and using $p/2>1$
preserves the bound.  Increasing $B$ absorbs the fixed logarithmic
power.  All ensuing coarse-bin witnesses are allowed to be anywhere
within $O(R^4)$ of their tested points.

\subsection{Flattened packets and their two energy budgets}

For a fine bin $\theta$ centered at $v$, use the original-axis map
$F_\theta=A_{h,v}$.  After removing its carrier, the fine-bin field is
\[
 h^{-n/2}U'_\theta(F_\theta z),
\]
where $U'_\theta$ is a child array of size $b_1$.  Partition its native
spacetime into smooth compact unit cells $T$, with a fixed enlargement
$T^*$ containing each cutoff support.  The resulting packets have
carrier $v$, width $h$ and duration $h^2$.  Their sizes and associated
densities can be taken as
\begin{align}
 s_{\theta,T}&=Ch^{-n/2}\max_{|\alpha|\leq N_*}
                     \sup_{T^*}|\partial^\alpha U'_\theta|,
                     \label{w:eq-flatsize}\\
 d_{\theta,T}&=(mh^{n+1})^{-1}
  \int(1+\dist(F_\theta z,T))^{-P_*}\,\dd\mu(z).
                     \label{w:eq-flatdensity}
\end{align}
Here $P_*,N_*$ are the fixed orders committed to in step
\textup{(ii)} of \Cref{w:choices}.  The envelope order $P=P_*$ was
chosen strictly above both $M+n+2$ and the lower envelope order in
\Cref{f:frame}.  Use it both for the full square-envelope weight
$w_T$ and for the density kernel.  Thus
\[
 \int w_T(z)\,\dd\mu(z)=mh^{n+1}d_{\theta,T},
\]
before replacing densities by their dyadic band values.  The order
$N_*$ includes the frame derivative order and the fixed reserves for the
normal Fourier expansion and cell Sobolev estimates.  Dividing each
localized profile by $s_{\theta,T}$ then gives the same fixed
low-order profile bound in every frame application.  Extra original
regularity needed by the sparse estimate does not redefine these
sizes or frame inputs.

Keep the cells needed on the physical domain and its macro
enlargements.  The pushed measure has ball parameter $Cmh^{n+1}$, so
$d_{\theta,T}\lesssim1$.

\begin{lemma}[Energy budgets after flattening]\label{w:flatten}
For each physical time bin of length $h^2$, including its fixed
enlargement,
\begin{equation}\label{w:eq-flat-unweighted}
 \sum_{\mathrm{bin}}s_{\theta,T}^2\lesssim h^{-n}.
\end{equation}
There is also the global bound, summed over all clusters, fine bins,
cells and time bins with bounded overlap,
\begin{equation}\label{w:eq-flat-weighted}
 \sum_{\theta,T}d_{\theta,T}^{\beta}s_{\theta,T}^2
 \lesssim_B h^{-n}b_1^{2n/(n+1)}D^{-c_s}.
\end{equation}
The first assertion is only a per-time-bin statement.
\end{lemma}

\begin{proof}
A physical time interval of length $h^2$ has bounded length in the
child time variable, even though a full child packet has duration
$b_1^2$.  At each such time, the Gram matrix of child atoms, also after
the fixed derivatives in \eqref{w:eq-flatsize}, has summable rows.
Integration by parts gives decay in the scaled carrier separation
$b_1(\omega-\omega')$, while compact spatial profiles localize
intercepts.  Transported phase-space counting bounds the number of
labels at each scaled distance.  All child frequencies are in a fixed
bounded box, so differentiation introduces no positive power of $b_1$.
Consequently
\[
 \sum_\theta\norm{\partial^\alpha U'_\theta(\cdot,t)}_2^2
 \lesssim\sum_\theta E_\theta\leq1
\]
for the required derivatives.  Unit-cell Sobolev, summed in space and
over the bounded child-time interval, proves
\eqref{w:eq-flat-unweighted}.

For the weighted bound apply \Cref{s:sparse}, including its fixed
derivative versions, to every full child array.  The pushforward of
$\mu$ has ball parameter $Cmh^{n+1}$ and its packet moment at scale
$b_1$ retains the same $D$, since
$h^{n+1}b_1^{n+1}=L^{n+1}$.  The choice
$b_1\asymp(BD)^{C_{10}}$ ensures that $D$ is at least a fixed positive
power of $b_1$ once $D$ is larger than a constant depending on $B$;
the remaining bounded $b_1,D$ cases cost only a constant depending on
$B,C_{10}$.  Also $D\leq b_1$ by the choice of $C_{10}$.

Convolve the pushed measure by the isotropic polynomial kernel used
in \eqref{w:eq-flatdensity}, and restrict to the suitably enlarged
time range.  Its ball and packet-moment bounds are preserved: translation
of a packet weight costs at most its fixed polynomial order, and
$P_*>M+n+2$ makes that cost integrable.  The convolved measure has
mass at least a constant times $mh^{n+1}d_{\theta,T}$ on the cells
where the suprema are taken.  The sparse estimate therefore gives
\eqref{w:eq-flat-weighted}, after division by the normalization in
\eqref{w:eq-flatsize}.  Cluster energies are disjoint and enlarged
physical time bins have bounded overlap, so this is a global sum.
\end{proof}

\subsection{Density bands and the frame estimate}

The two budgets will serve different parts of the summation.  Very
small densities will pay the polynomial number of clusters and time
bins, so their energies can be bounded one time bin at a time.  For
the remaining densities we will choose transverse coarse bins from
one common density band.  The corresponding density weight can then
remain attached to each packet energy and be summed with the global
weighted budget.  We now construct these two kinds of tuples and verify the
hypotheses of the frame estimate.

Write the densities in dyadic bands $d_\ell=C_d2^{-\ell}$,
$\ell\geq0$, with fixed upper constant $C_d$.  Only $O(\log h)$
indices occur on the needed cells.  Indeed their distances from the
measure support are bounded by a fixed power of $h$, while
\eqref{w:eq-incidence-mass} gives a polynomial lower bound on total
mass; the strictly positive kernel in \eqref{w:eq-flatdensity} then
gives $d_{\theta,T}\geq h^{-C}$.

At a witness for each significant coarse bin, at least one band field
has modulus at least
\begin{equation}\label{w:eq-bandlevel}
 A_\ell=c'R^{-n}\lambda(1+\ell)^{-2}.
\end{equation}
This follows by threshold splitting, since
$\sum_{\ell\geq0}(1+\ell)^{-2}<\infty$.  Choose one such band for
each significant bin.  There are two possibilities:
\begin{enumerate}[label=\textup{(\roman*)}]
\item A chosen band is tiny, $d_\ell<X^{-C_{11}}$.  Starting from its
center, complete it to a robust $n=k+1$ tuple of significant coarse
bins in tangential frequency space.
\item All chosen bands are non-tiny.  There are only $O(C_{11}\log X)$
colors.  Some single color supplies a robust $n$-tuple.
\end{enumerate}
Both statements follow from the strip-union exclusion: successive
maximal affine heights can all be chosen at least $C_*/R$.  If a color
failed to supply such a simplex, it would lie in one strip.  The
determinant lower bound for a successful tuple is $cR^{-k}$ throughout
the bins when $C_*$ is large enough.

Fix a cluster, a time bin of length $h^2$, an ordered coarse tuple and
its band indices $\ell_i$, $1\leq i\leq n$.  Let $M_0$ be the mass
of tested points assigned to this tuple, and let
$E_i=\sum s_{\theta,T}^2$ over its $i$th part, with the fixed time
enlargement.  A tuple of mass at most $mh^{2n-c_4}$ can be discarded,
where $c_4>0$ is a small fixed dimensional constant.  The total number
of tuples is $X^C(\log h)^C$, so their relative discarded mass, by
\eqref{w:eq-incidence-mass}, is at most
\[
 X^C(\log h)^C(h/L)^{2n}h^{-c_4}=o(1).
\]
This is one of the requirements imposed on the later $\kappa$.

For every remaining tuple, Markov and \eqref{w:eq-flatdensity} retain
at least half its mass where all full square envelopes are bounded by
\begin{equation}\label{w:eq-Bi}
 B_i=C\left(\frac{mh^{n+1}d_{\ell_i}E_i}{M_0}\right)^{1/2}.
\end{equation}
They remain comparable at every point of the macro containing the
tested point, since $R^4\ll h$.  Thus they also bound the envelopes
at the separate coarse-bin witnesses.  Zero-energy cases need no
consideration.  By \eqref{w:eq-flat-unweighted} and the mass cutoff,
\begin{equation}\label{w:eq-framecutoff}
 B_i\lesssim h^{(1-2n+c_4)/2},\qquad A_{\ell_i}/B_i\geq h^{c_3}
\end{equation}
for a fixed $c_3>0$.  For the latter assertion, use
$g\geq D^{-1}$, $\chi>L^{-v_0}$ and the margin
\eqref{w:eq-threshold-margin}.  Choose $c_4,v_0,c_3$ as sufficiently
small dimensional fractions of that margin; the later choice of
$\kappa$ pays every fixed $X$ power in $L/h$, $R$ and the threshold,
and logarithmic factors are smaller than any fixed positive power of
$h$.  This verifies the frame cutoff without allowing a geometric
power to determine $u$ retroactively.

\begin{lemma}[A tuple bound]\label{w:tuple}
For each retained tuple,
\begin{equation}\label{w:eq-tuple}
 \lambda^2 M_0\leq CR^C(1+\textstyle\sum_i\ell_i)^C
 m\chi^{1-\beta}h^{n+2n/(n+1)}
 \left(\prod_{i=1}^n d_{\ell_i}E_i\right)^{\beta/n}
 \left(\sum_{i=1}^nE_i\right)^{1-\beta}.
\end{equation}
The $R$ exponent is independent of the number of plate clusters and
time bins and of $C_{10},C_{11}$.
\end{lemma}

\begin{proof}
Partition the sheared normal spatial coordinate into slabs of width
$h$, with bounded enlargements to accommodate witnesses.  The projected
measure in tangential spacetime has unit-cube mass at most $Cm\chi h$
on each slab, by unit covering in the normal coordinate.

Only nearby normal and time labels of compact packets contribute in
a fixed slab and time bin.  Recenter at the time-bin origin.  The
residual normal frequency is $O(1/h)$, so its drift during $h^2$ is
$O(h)$.  Projected tangential packets consequently have bounded
phase-space multiplicity.  Rotating the lattice changes only bounded
multiplicities and comparable polynomial weights.  The omitted normal
and time indices have bounded normalized displacements, so the projected
envelopes are bounded by a constant times the full envelopes.

In the normal coordinate divided by $h$, expand each cutoff profile
in a Fourier series.  The residual normal carrier and time phase have
bounded normalized derivatives and can be included in that profile.
The exterior Fourier phase has modulus one.  For mode $m\in\Z$,
the tangential packet amplitude and envelope gain
$\rho_m\lesssim(1+|m|)^{-A_*}$, with any required fixed $A_*$ supplied
by the original derivative order.  Split thresholds by
$\gamma_m=c(1+|m|)^{-2}$.  The frame cutoff improves rather than
deteriorates under the factor $\gamma_m/\rho_m$.

Witnesses need not be on the tested point's projected unit cube.
Translate the tangential spacetime packets by an integer offset,
separately for each family, so that its witness lies in that cube.
Each offset is $O(R^4)$, so summing all choices costs $R^C$.
Counting, transversality and envelope comparability are preserved.

Apply \Cref{f:frame}, normalizing family $i$ by $B_i$, and multiply
the resulting cube count by $Cm\chi h$.  Before summing modes, the
right side contains the decay factors
\[
 \prod_{i=1}^n(\rho_{m_i}/\gamma_{m_i})^{4/(kn)}
 \quad\hbox{and, in its $j$th summand,}
 \quad(\rho_{m_j}/\gamma_{m_j})^2.
\]
Every family has a positive decay exponent, even when it does not
index the energy summand.  All independent mode sums therefore
converge for fixed large $A_*$.  Sum slabs at this stage, keeping the
same global $B_i$ from \eqref{w:eq-Bi} in every slab and putting only
its nearby packet energies in the linear sum.  Each packet occurs in
boundedly many enlarged slabs.  We obtain
\begin{equation}\label{w:eq-frame-before-solve}
 M_0\leq CR^C m\chi h
 \prod_{i=1}^n(B_i/A_{\ell_i})^{4/(kn)}
 \sum_{i=1}^n\frac{h^{k+2}E_i}{A_{\ell_i}^2}.
\end{equation}
Summing slabs after solving a nonlinear mass inequality instead would
lose a multiplicity factor; \eqref{w:eq-frame-before-solve} avoids it.

Substitute \eqref{w:eq-Bi} and \eqref{w:eq-bandlevel}.  Up to the
displayed $R$ and band-index factors, the result is
\[
 M_0^{1+2/k}\leq C R^C(1+\textstyle\sum_i\ell_i)^C
 m^{1+2/k}\chi
 h^{n+2+2(n+1)/k}\lambda^{-2-4/k}
 \left(\prod_i d_{\ell_i}E_i\right)^{2/(kn)}\sum_iE_i.
\]
Raise both sides to $k/(k+2)=1-\beta$.  The exponent of $h$ becomes
$n+2n/(n+1)$ and the product exponent becomes $\beta/n$; the
$\lambda$ exponent is exactly $-2$.  This is
\eqref{w:eq-tuple}.
\end{proof}

\subsection{Summation and completion of the induction}

\begin{proof}[Completion of the proof of \Cref{w:weak}]
First sum the tuples with a tiny band.  The per-bin bound
$E_i\lesssim h^{-n}$ gives
\[
 \left(\prod_i E_i\right)^{\beta/n}
 \left(\sum_iE_i\right)^{1-\beta}\lesssim h^{-n}.
\]
Here the polynomial number of clusters and the
$O((L/h)^2)=X^{2C_{10}}$ time bins may be counted explicitly.  Their
total power of $X$ is fixed before $C_{11}$.  To sum the density bands,
extend the range to all nonnegative indices and use
\begin{equation}\label{w:eq-tinysum}
 \sum_{\min_i d_{\ell_i}<X^{-C_{11}}}
 (1+\textstyle\sum_i\ell_i)^C\prod_i d_{\ell_i}^{\beta/n}
 \leq C X^{-\beta C_{11}/(2n)}.
\end{equation}
Indeed assign a tiny index, bound the polynomial by a product of
polynomials, and absorb each into half its exponential decay.  The
remaining series are geometric.  There is no factor $\log h$ in
\eqref{w:eq-tinysum}.  Choose $C_{11}$ large enough to pay all the
previous cluster, time-bin and tuple powers and the desired negative
power of $D$.

For a common non-tiny band, the energy expression in
\eqref{w:eq-tuple} is bounded by
\[
 C d_\ell^\beta\sum_i E_i.
\]
Each packet has one density index, and can occur in only $R^C$ coarse
tuple and witness-offset choices.  Normal slabs and enlarged time
bins have bounded overlap.  Cluster labels are disjoint coefficient
subarrays.  Hence this sum is charged directly to the global weighted
bound \eqref{w:eq-flat-weighted}, with no factor counting clusters or
time bins.  Also
$\ell\leq C+C_{11}\log_2X$, so its band factor is only
$C(C_{11})(1+\log X)^C$.  The scale factors cancel as
\[
 h^{n+2n/(n+1)}h^{-n}b_1^{2n/(n+1)}
 =L^{2n/(n+1)}.
\]
The choices of $u$ and then $B$ absorb $R^{C_f}$ and the fixed
logarithmic powers into part of the sparse gain.  Combining both
types of tuples, and recalling that a fixed fraction of $M_Y$ remained,
gives
\begin{equation}\label{w:eq-finalL2}
 \lambda^2M_Y\leq C_Bm\chi^{1-\beta}
                   L^{2n/(n+1)}D^{-c_s/2}.
\end{equation}

We now verify the promised polynomial level bound, before choosing
$p-2$.  From \eqref{w:eq-incidence-mass},
$M_Y\geq cmL^{2n}X^{-C_{\mathrm{geo}}}$.  Substitute this and
$\lambda=\chi^agL^{-e}$ in \eqref{w:eq-finalL2}.  Since
$2n-2e=2n/(n+1)$, we obtain
\[
 g^2\leq C_B\chi^{\zeta}X^{C_{\mathrm{geo}}}D^{-c_s/2}
 \leq C_BX^{C_{\mathrm{geo}}}.
\]
Thus $g\leq C_BX^{C_g}$ with $C_g$ fixed by geometry, independently
of $\kappa,p-2,C_{\mathrm{ind}}$.  This confirms the noncircular order
of choices in \Cref{w:choices}.

Finally \eqref{w:eq-finalL2} and $\chi^{\zeta}\leq1$ imply
\[
 M_Y\leq C_BmL^{2n}g^{-2}D^{-c_s/2}
 \leq C'_BmL^{2n}g^{-p}D^{-\sigma},
\]
provided
\[
 C_g(p-2)+\sigma<c_s/2.
\]
This requirement is compatible with the easy-range inequalities,
since $p-2$ is chosen after all geometric exponents and after $\kappa$.
All uses of induction removed small fractions of the hypothetical
mass: their ratios were negative powers of $X$, or
$R^{-c_2}D^\sigma$, fixed before choosing $C_{\mathrm{ind}}$.
Every remaining constant $C'_B$ is independent of that induction
constant.  Taking $C_{\mathrm{ind}}>C'_B$ contradicts
\eqref{w:eq-counterlevel}.  Both inductions close, proving
\Cref{w:weak}.
\end{proof}

